\chapter{Contextualización y Antecedentes}\label{cap:contexto}
La inteligencia artificial (IA) ha pasado de ser un campo de investigación especializado a constituir una tecnología transversal presente en la salud, la industria, la educación y los servicios públicos. Este capítulo sitúa el proyecto mediante un enfoque de embudo que recorre cinco niveles de análisis: internacional, nacional, departamental, regional/provincial y local, hasta llegar a la realidad concreta de los estudiantes de Programación 1 de la FICCT.

\section{Nivel internacional}
Las redes neuronales artificiales tienen una historia que se remonta al modelo lógico de neurona de \citet{mcculloch1943} y al perceptrón de \citet{rosenblatt1958}. El entusiasmo inicial se frenó cuando \citet{minsky1969} demostraron formalmente que un perceptrón de una sola capa no puede resolver problemas no linealmente separables como la compuerta XOR. El campo resurgió con la difusión del algoritmo de retropropagación del error \citep{rumelhart1986} y con los resultados teóricos que establecen que una red con al menos una capa oculta es un aproximador universal de funciones continuas \citep{cybenko1989,hornik1989}. Desde la década de 2010, el aprendizaje profundo ha transformado la visión por computador, el procesamiento del lenguaje y numerosos dominios científicos \citep{lecun2015,goodfellow2016}.

El informe \emph{AI Index} de la Universidad de Stanford documenta el crecimiento sostenido de la inversión, las publicaciones y los modelos de IA a escala mundial, así como la concentración de su desarrollo en un número reducido de organizaciones con gran capacidad de cómputo \citep{maslej2024}. En paralelo, la UNESCO ha subrayado la necesidad de formar ciudadanos y profesionales capaces de comprender críticamente los sistemas de IA, y no únicamente de utilizarlos \citep{unesco2022}. Esta tensión revela una brecha formativa: una gran cantidad de usuarios emplea \emph{frameworks} de alto nivel como cajas negras, mientras que comparativamente pocos desarrolladores dominan los fundamentos matemáticos ---álgebra lineal, cálculo diferencial y optimización--- que explican por qué y cómo aprende una red neuronal \citep{deisenroth2020}.

En el ámbito de la ingeniería de sistemas, el lenguaje Rust se ha consolidado como una alternativa que combina el rendimiento de C y C++ con garantías de seguridad de memoria verificadas en tiempo de compilación \citep{matsakis2014,jung2018}. La encuesta anual de desarrolladores de Stack Overflow lo ha ubicado de forma reiterada como el lenguaje más «admirado» por quienes lo utilizan \citep{stackoverflow2024}. Estas características lo vuelven atractivo para la construcción de componentes de alto rendimiento, incluidos los motores numéricos sobre los que se apoya el aprendizaje automático.

\section{Nivel nacional (Bolivia)}
En Bolivia, el marco normativo de las tecnologías de información y comunicación se sustenta en la Ley N.º 164, Ley General de Telecomunicaciones, Tecnologías de Información y Comunicación, que reconoce el acceso a las TIC como un servicio de interés público y promueve el uso de software libre y estándares abiertos en el Estado \citep{bolivia2011}. Este marco favorece iniciativas académicas que, como la presente, se basan en herramientas abiertas y gratuitas.

La transformación digital del país plantea una demanda creciente de profesionales capaces no solo de consumir tecnología, sino de comprenderla y adaptarla a las necesidades locales. En el campo específico de la IA, la formación universitaria todavía tiende a priorizar el uso de herramientas sobre la comprensión de sus fundamentos. Un proyecto que reconstruye una red neuronal desde sus bases matemáticas contribuye a cerrar esa brecha y a formar talento con capacidad de innovación propia.

\begin{nota}[Dato pendiente de verificación]
Para la versión final se recomienda incorporar indicadores nacionales verificables (por ejemplo, estadísticas oficiales de matrícula en carreras de computación o de conectividad) citados en formato APA 7 desde fuentes institucionales. En este borrador no se incluyen cifras que no hayan sido verificadas.
\end{nota}

\section{Nivel departamental (Santa Cruz)}
Santa Cruz es el departamento con mayor dinamismo económico del país y concentra una parte importante de la actividad empresarial y de servicios. Su ecosistema tecnológico incluye empresas de desarrollo de software, emprendimientos digitales y comunidades técnicas que organizan eventos, talleres y encuentros de programación. En este entorno, la demanda de ingenieros con una base matemática y algorítmica sólida es cada vez mayor, porque los productos basados en datos e IA requieren profesionales capaces de diagnosticar, ajustar y optimizar modelos en lugar de limitarse a invocar bibliotecas.

\section{Nivel regional y provincial}
El acceso al conocimiento tecnológico avanzado tiende a concentrarse en las capitales. Para los estudiantes de las provincias del departamento, la dependencia de hardware costoso, de servicios de cómputo en la nube o de licencias pagadas constituye una barrera concreta. El enfoque de este proyecto ---una red neuronal escrita en Rust puro, sin dependencias externas, que se compila en un binario de aproximadamente \BinKiB{}~KiB y entrena los tres experimentos en fracciones de segundo en un computador convencional--- demuestra que es posible estudiar los fundamentos del aprendizaje automático con recursos modestos y materiales de código abierto, lo que favorece la descentralización del conocimiento.

\section{Nivel local (UAGRM -- FICCT)}
En la FICCT, la asignatura Programación 1 (\Sigla) introduce a los estudiantes en el pensamiento algorítmico, las estructuras de control y la modularidad del código. Con frecuencia, los estudiantes de primeros semestres perciben la IA como un tema lejano y reservado a especialistas. Los Semilleros de Investigación impulsados por el IICCT y la Feria Facultativa de Proyectos ofrecen un espacio para acercar estos temas a etapas tempranas de la formación.

Este proyecto aprovecha ese espacio para demostrar que los conceptos de Programación 1 ---funciones, estructuras, ciclos, arreglos y pruebas--- son suficientes para construir, paso a paso, un sistema que aprende a partir de datos. De este modo, el trabajo se convierte en un recurso educativo reutilizable por otros estudiantes de la facultad.

\section{Antecedentes}
Entre los antecedentes conceptuales del proyecto destacan los textos de \citet{haykin2009}, \citet{bishop2006} y \citet{goodfellow2016}, que desarrollan con rigor la teoría de redes neuronales, y el libro abierto de \citet{nielsen2015}, que construye una red neuronal desde cero con fines didácticos. En el ámbito práctico, \citet{lecun1998} sistematizaron recomendaciones para entrenar redes mediante retropropagación, mientras que \citet{glorot2010} y \citet{he2015} propusieron esquemas de inicialización de pesos que hoy son estándar. Respecto a Rust, la documentación oficial \citep{klabnik2023} y el texto de \citet{blandy2021} constituyen las referencias principales sobre el lenguaje y su modelo de propiedad (\emph{ownership}).
